\title[TITLETOKEN]{TITLETOKEN}
\subtitle{Laboratory Presentation · \labversionlabel}
\author{}
\institute{Institution Name}
\date{DATEPLACEHOLDER}

\begin{document}

\begin{frame}[plain]
  \titlepage
\end{frame}

\begin{frame}{Laboratory Workflow}
  \slidetag{WORKFLOW}
  \vspace{0.7em}
  \begin{columns}[T,onlytextwidth]
    \begin{column}{0.48\textwidth}
      \resultcard{Step 1}{Understand the physical model}
      \vspace{0.6em}
      \resultcard{Step 2}{Assemble the apparatus and measure}
    \end{column}
    \begin{column}{0.48\textwidth}
      \resultcard{Step 3}{Process the results}
      \vspace{0.6em}
      \resultcard{Step 4}{State the conclusion}
    \end{column}
  \end{columns}
\end{frame}

\begin{frame}{Objective and Measurable Outcome}
  \begin{columns}[T,onlytextwidth]
    \begin{column}{0.56\textwidth}
      \begin{block}{Objective}
        State a measurable objective and identify the physical quantity to
        be determined.
      \end{block}
    \end{column}
    \begin{column}{0.39\textwidth}
      \resultcard{Deliverable}{Plot, calculation, and conclusion}
    \end{column}
  \end{columns}
\end{frame}

\begin{frame}{Equipment and Safety}
  \begin{columns}[T,onlytextwidth]
    \begin{column}{0.48\textwidth}
      \slidetag{EQUIPMENT}
      \vspace{0.6em}
      \begin{itemize}
        \item list the instruments;
        \item specify their measurement ranges;
        \item note relevant specifications.
      \end{itemize}
    \end{column}
    \begin{column}{0.48\textwidth}
      \begin{alertblock}{Safety}
        Include only the instructor's safety requirements that apply to this apparatus.
      \end{alertblock}
    \end{column}
  \end{columns}
\end{frame}

\begin{frame}{Physical Model}
  \keyformula{$F=ma$}
  \vspace{0.8em}
  \begin{itemize}
    \item Define the notation before the first equation.
    \item Explain the physical meaning of the model.
    \item State its assumptions and range of validity.
  \end{itemize}
\end{frame}

\begin{frame}{Procedure}
  \begin{enumerate}
    \item Assemble and check the apparatus.
    \item Perform a series of measurements.
    \item Record the raw data in a table.
    \item Process the results and estimate the uncertainty.
    \item State a conclusion that addresses the objective.
  \end{enumerate}
\end{frame}

\begin{frame}{Check-for-Understanding Question}
  \slidetag{SELF-CHECK}
  \vspace{0.7em}
  \begin{block}{Question}
    Add a question that assesses understanding of the laboratory topic.
  \end{block}
\end{frame}

\begin{solutionframe}[Check-for-Understanding Solution]
  \begin{exampleblock}{Solution}
    Give a complete, explained answer rather than only the final equation.
  \end{exampleblock}
\end{solutionframe}

\begin{frame}{Report Requirements}
  \begin{columns}[T,onlytextwidth]
    \begin{column}{0.48\textwidth}
      \begin{itemize}
        \item objective and apparatus diagram;
        \item raw-data table;
        \item sample calculation.
      \end{itemize}
    \end{column}
    \begin{column}{0.48\textwidth}
      \begin{itemize}
        \item plot or fitted model;
        \item uncertainty estimate;
        \item physics-based conclusion.
      \end{itemize}
    \end{column}
  \end{columns}
\end{frame}

\end{document}
